\begin{ouverture}[geometrie]{16}{Parallélogrammes et aires}{Partie 3 · Espace et géométrie}
{Côtés · Diagonales · Rectangle · Losange · Carré}
\objectif{Reconnaître et construire des parallélogrammes.}
\objectif{Choisir une propriété caractéristique.}
\objectif{Calculer des aires simples et composées.}
\prerequis{Parallélisme, symétrie centrale et hauteurs.}
\end{ouverture}
\pagestyle{otMain}\essentiel\fontsize{11}{15}\selectfont
\phantomsection\label{cours-16-1}
\essSection{1}{Définition et propriétés caractéristiques}
\begin{definition}Un parallélogramme est un quadrilatère dont les côtés opposés sont parallèles. Ses côtés opposés sont aussi de même longueur ; ses diagonales ont le même milieu. Leur milieu commun est un centre de symétrie.\end{definition}
\begin{definition}Un quadrilatère a quatre côtés, un pentagone cinq, un hexagone six. Un trapèze possède au moins une paire de côtés opposés parallèles ; certains usages réservent ce mot au cas d’une seule paire. Préciser la convention quand elle intervient.\end{definition}
\label{fig-cours-16-1-cours-parallelogramme}\figureGeo{cours-parallelogramme}

\begin{propriete}Pour un quadrilatère non croisé, chacune de ces conditions permet de prouver qu’il est un parallélogramme : les côtés opposés sont parallèles deux à deux ; les côtés opposés sont égaux deux à deux ; les diagonales ont le même milieu ; deux côtés opposés sont à la fois parallèles et de même longueur. Par exemple, si [AC] et [BD] ont le même milieu O, le demi-tour de centre O échange A et C, B et D : (AB) a pour image (CD), qui lui est donc parallèle, et de même pour (BC) et (DA).\end{propriete}
\begin{savoirfaire}[Construire avec côtés ou diagonales]Pour construire ABCD à partir de A, B et C consécutifs, tracer par A la parallèle à (BC), puis par C la parallèle à (AB). Leur intersection est D. On peut aussi construire le milieu O de [AC] puis le symétrique de B par rapport à O.\end{savoirfaire}
\clearpage
\phantomsection\label{cours-16-2}
\essSection{2}{Cas particuliers et aires}
\begin{definition}Un rectangle est un parallélogramme avec un angle droit. Un losange est un parallélogramme avec deux côtés consécutifs égaux. Un carré est à la fois rectangle et losange.\end{definition}
\begin{propriete}Pour un parallélogramme : des diagonales de même longueur caractérisent un rectangle ; des diagonales perpendiculaires caractérisent un losange. Si les deux conditions sont réunies, c’est un carré. Sans l’hypothèse « parallélogramme », ces tests ne suffisent pas.\end{propriete}
\label{fig-cours-16-2-cours-aire-parallelogramme}\figureGeo{cours-aire-parallelogramme}

\begin{propriete}L’aire d’un parallélogramme est $b\times h$, avec une base et la hauteur perpendiculaire correspondante. Elle s’explique par découpage et déplacement d’un triangle pour former un rectangle.\end{propriete}
\begin{savoirfaire}[Convertir avant de combiner]Pour une figure composée, découper en figures dont on sait calculer l’aire, ou retirer une partie d’une figure simple. $1$ m$^2=100$ dm$^2=10000$ cm$^2$.\end{savoirfaire}

\clearpage\rubrique{Découverte}
\begin{blocLivre}{ex-16-1}\begin{exercice}[1]{Côtés opposés}
ABCD est un parallélogramme, $AB=7$ cm et $BC=4$ cm. Donner CD et AD.
\label{fig-ex-16-1}\figureGeo{para-1}
\end{exercice}\end{blocLivre}
\begin{blocLivre}{ex-16-2}\begin{exercice}[1]{Diagonales et milieu}
Dans le parallélogramme ABCD, les diagonales se coupent en O. Si $AC=10$ cm et $BD=6$ cm, donner AO et BO.
\label{fig-ex-16-2}\figureGeo{para-2}
\end{exercice}\end{blocLivre}
\begin{blocLivre}{ex-16-3}\begin{exercice}[1]{Un rectangle}
ABCD est un parallélogramme et $\widehat A=90^\circ$. Préciser sa nature.
\label{fig-ex-16-3}\figureGeo{para-3}
\end{exercice}\end{blocLivre}
\begin{blocLivre}{ex-16-4}\begin{exercice}[1]{Aire directe}
Un parallélogramme a une base de 8 cm et la hauteur correspondante de 3 cm. Calculer son aire.
\label{fig-ex-16-4}\figureGeo{para-4}
\end{exercice}\end{blocLivre}
\begin{blocLivre}{ex-16-5}\begin{exercice}[1]{Un losange}
ABCD est un parallélogramme tel que $AB=BC=5$ cm. Quelle est sa nature ?
\label{fig-ex-16-5}\figureGeo{para-5}
\end{exercice}\end{blocLivre}

\rubrique{Entraînement}
\begin{blocLivre}{ex-16-6}\begin{exercice}[2]{Construire le quatrième sommet}
Construire D pour que A, B et C soient trois sommets consécutifs d’un parallélogramme ABCD.
\label{fig-ex-16-6}\figureGeo{para-6}
\end{exercice}\end{blocLivre}
\begin{blocLivre}{ex-16-7}\begin{exercice}[2]{Construire par les diagonales}
A, B et C sont donnés. Construire D afin que ABCD soit un parallélogramme, en utilisant le milieu de [AC].
\label{fig-ex-16-7}\figureGeo{para-7}
\end{exercice}\end{blocLivre}
\begin{blocLivre}{ex-16-8}\begin{exercice}[2]{Rectangle par diagonales}
Les diagonales d’un quadrilatère ont le même milieu et mesurent chacune 8 cm. Quelle nature peut-on prouver ?
\label{fig-ex-16-8}\figureGeo{para-8}
\end{exercice}\end{blocLivre}
\begin{blocLivre}{ex-16-9}\begin{exercice}[2]{Losange par diagonales}
Les diagonales d’un quadrilatère se coupent perpendiculairement en leur milieu. Quelle est sa nature ?
\label{fig-ex-16-9}\figureGeo{para-9}
\end{exercice}\end{blocLivre}
\begin{blocLivre}{ex-16-10}\begin{exercice}[2]{Un carré à reconnaître}
Un parallélogramme a des diagonales égales et perpendiculaires. Justifier que c’est un carré.
\label{fig-ex-16-10}\figureGeo{para-10}
\end{exercice}\end{blocLivre}

\rubrique{Problèmes}
\begin{blocLivre}{ex-16-11}\begin{exercice}[3]{Hauteur inconnue}
Un parallélogramme de base 12 cm a une aire de 54 cm$^2$. Trouver sa hauteur.
\label{fig-ex-16-11}\figureGeo{para-11}
\end{exercice}\end{blocLivre}
\begin{blocLivre}{ex-16-12}\begin{exercice}[3]{Unités différentes}
La base d’un parallélogramme vaut 1,2 m et sa hauteur 35 cm. Calculer son aire en m$^2$ puis en cm$^2$.
\label{fig-ex-16-12}\figureGeo{para-12}
\end{exercice}\end{blocLivre}
\begin{blocLivre}{ex-16-13}\begin{exercice}[3]{Une figure composée}
Une figure est formée d’un rectangle de 8 cm sur 4 cm et d’un triangle accolé sur un côté de 8 cm, de hauteur 3 cm. Calculer l’aire totale.
\label{fig-ex-16-13}\figureGeo{para-13}
\end{exercice}\end{blocLivre}
\begin{blocLivre}{ex-16-14}\begin{exercice}[3]{Une hypothèse manquante}
Une élève affirme : « Si les diagonales sont perpendiculaires, le quadrilatère est un losange. » Examiner le contre-exemple dessiné.
\label{fig-ex-16-14}\figureGeo{para-14}
\end{exercice}\end{blocLivre}
\begin{blocLivre}{ex-16-15}\begin{exercice}[3]{Deux méthodes pour une aire}
Un losange a des diagonales de 10 cm et 6 cm. Les diagonales sont perpendiculaires et se coupent en leur milieu. Calculer son aire par découpage en quatre triangles.
\label{fig-ex-16-15}\figureGeo{para-15}
\end{exercice}\end{blocLivre}

\clearpage
